% Lions-Style Commentary on SM89 SASS Instructions
%
% Following John Lions' "Commentary on UNIX 6th Edition" (1977):
% numbered paragraphs, each explaining one aspect of the machine,
% cross-referenced to source (probe files) and raw data (SASS disassembly).
%
% Each paragraph number (\paragraph) corresponds to a row in the
% multidimensional table (Table~\ref{tab:multidim}).

\section{A Commentary on the SM89 Instruction Set}

\noindent\textit{%
What follows is a guided reading of the Ada Lovelace SM89 instruction set,
in the style of Lions' Commentary on UNIX.  Each numbered paragraph explains
one instruction class, what it reveals about the microarchitecture, and how
we measured it.  The reader should have Table~\ref{tab:multidim} and the
pipeline diagram (Figure~\ref{fig:pipeline}) open alongside.}

\bigskip

%----------------------------------------------------------------------
\subsection{Integer Arithmetic Pipeline}
\label{par:int32}

\paragraph{2.1}
\texttt{IADD3} is the fundamental three-input integer addition.  Its measured
latency of 2.51 cycles and throughput of 68.2 ops/clk/SM tell us two things.
First, the INT32 pipeline is the fastest execution unit on SM89 -- faster than
the FP32 FMA pipe by nearly 2x in latency.  Second, the three-input form is
native to the hardware, not a fused pair of two-input adds; the compiler uses
\texttt{IADD3} with the third operand set to zero when only two inputs are
needed, confirming that this is the canonical integer add.

\textit{Source:} \texttt{probes/probe\_int\_arith.cu}, measurement chain depth 512.

\paragraph{2.2}
\texttt{IMAD} (integer multiply-add) shares the 4.54-cycle latency of
\texttt{FFMA}, suggesting the INT32 multiplier may share datapath resources
with the FP32 FMA unit, or at minimum occupies the same pipeline depth.  The
compiler uses \texttt{IMAD.IADD} as a synonym for \texttt{IADD3} in certain
contexts, and \texttt{IMAD.X} for carry-propagating multiword arithmetic.
The \texttt{.IADD} and \texttt{.X} suffixes are not separate instructions but
modifier bits in the same opcode encoding.

\textit{Source:} \texttt{probes/probe\_int\_arith.cu}, encoding analysis in
\texttt{results/20260306\_190541/ENCODING\_ANALYSIS.md}.

%----------------------------------------------------------------------
\subsection{Floating-Point Arithmetic Pipeline}
\label{par:fp32}

\paragraph{3.1}
\texttt{FFMA} at 4.54 cycles is the workhorse of GPU compute.  Every matrix
multiplication, convolution, and most shader computations reduce to long
chains of FFMA instructions.  The measured throughput of 44.6 ops/clk/SM,
combined with the 128 CUDA cores per SM on Ada, implies that each core can
sustain slightly more than one FFMA every 4 clocks -- consistent with the
4-cycle pipeline depth.  The MLP forward kernel in \texttt{instant\_ngp/}
demonstrates that 8-wide ILP across independent accumulators is needed to
saturate this pipeline; the compiler's default 1-wide chain leaves 7/8 of
the pipeline idle.

\textit{Source:} \texttt{microbench/microbench\_latency.cu} (chain depth 512),
cross-validated with ncu \texttt{smsp\_\_cycles\_elapsed} and
\texttt{sm\_\_throughput.avg.pct\_of\_peak\_sustained\_elapsed}.

\paragraph{3.2}
\texttt{FMNMX} (float min/max) shares the FMA pipeline at identical latency.
This is significant for activation functions: a ReLU implemented as
\texttt{FMNMX Rd, Ra, RZ, !PT} (take max of value and zero) costs exactly
one pipeline slot, the same as an FFMA.  The hand-optimized \texttt{mlp\_forward}
kernel exploits this by replacing a branch-based ReLU with a single FMNMX,
eliminating a branch and a comparison.

%----------------------------------------------------------------------
\subsection{Half-Precision and BF16 Pipeline}
\label{par:fp16}

\paragraph{4.1}
\texttt{HFMA2} at 4.01 cycles is the lowest-latency FMA operation measured
on SM89 -- faster than both FP32 \texttt{FFMA} and INT32 \texttt{IMAD}.
The \texttt{2} suffix indicates SIMD-2 operation: two FP16 values packed
into a single 32-bit register, processed in a single cycle.  This means
the effective throughput for FP16 is 2x the FP32 throughput at lower latency.
BF16 (\texttt{bf162}) operations achieve identical latency, confirming that
the BF16 datapath shares the FP16 SIMD pipeline with no additional penalty.

\textit{Source:} ncu SM\_cyc measurements across \texttt{microbench\_throughput.cu}
variants, cross-validated with \texttt{k\_hfma2\_bf16} showing the lowest
\texttt{SM\_cyc} of all measured kernels.

%----------------------------------------------------------------------
\subsection{FP64 Pipeline}
\label{par:fp64}

\paragraph{5.1}
\texttt{DFMA} at 54.48 cycles reveals the Ada gaming SKU's deliberate FP64
throttling.  The 1:64 FP64:FP32 throughput ratio (confirmed by
\texttt{sm\_\_throughput} measurements) means FP64 operations are
\emph{compute-bound} even on bandwidth-bound workloads -- the only precision
tier where this is true for scalar LBM kernels.  This has a direct
architectural implication: FP64 and double-double (DD) kernels cannot benefit
from the same memory optimization strategies that accelerate FP32/INT8 kernels.
The DD kernel's 2.85x slowdown vs FP32 is not a memory bottleneck but a raw
ALU throughput wall.

\textit{Source:} \texttt{probes/probe\_fp64\_arithmetic.cu},
\texttt{microbench/microbench\_latency\_corrected.cu}.

%----------------------------------------------------------------------
\subsection{Special Function Unit (SFU / MUFU)}
\label{par:sfu}

\paragraph{6.1}
The SFU executes transcendental approximations (\texttt{MUFU.EX2},
\texttt{MUFU.RCP}, \texttt{MUFU.RSQ}, \texttt{MUFU.SIN}, \texttt{MUFU.COS},
\texttt{MUFU.LG2}).  The wide latency spread -- 17.56 cycles for
\texttt{MUFU.EX2} vs 41.55 cycles for \texttt{MUFU.RCP} -- reveals that
these are not a single unit but a pipeline with operation-dependent depth.
Reciprocal (\texttt{RCP}) requires more Newton-Raphson iterations internally
than exp2 (\texttt{EX2}), which is a direct table lookup with interpolation.

\paragraph{6.2}
A critical finding for kernel optimization: \texttt{MUFU.RCP64H} (the FP64
reciprocal seed) shares the same ncu \texttt{SM\_cyc} profile as the BREV
tiling kernel, placing it in the SFU pipeline rather than the FP64 ALU.
This means FP64 division can partially overlap with FP32 computation if
the reciprocal seed is issued to the SFU while the FP64 ALU handles
multiply-accumulate.

\textit{Source:} \texttt{probes/probe\_mufu.cu},
\texttt{probes/probe\_transcendentals.cu}.

%----------------------------------------------------------------------
\subsection{Tensor Core Operations}
\label{par:tensor}

\paragraph{7.1}
\texttt{HMMA.1688.F32.TF32} at 66.19 measured cycles is the first published
tensor core latency for Ada Lovelace consumer hardware.  The \texttt{1688}
encoding specifies the M16-N8-K8 tile shape.  Unlike ALU operations that
execute per-thread, tensor core operations are warp-cooperative: all 32
threads in a warp contribute fragments to a single matrix multiply-accumulate.
The high absolute latency is offset by the massive data parallelism: each
invocation processes 16x8x8 = 1024 multiply-accumulate operations.

\textit{Source:} \texttt{probes/probe\_tensor\_all\_widths.cu},
\texttt{microbench/microbench\_latency\_tensor\_all.cu}.

%----------------------------------------------------------------------
\subsection{Bitwise and Logic Operations}
\label{par:bitwise}

\paragraph{8.1}
\texttt{LOP3.LUT} is the universal 3-input logic operation.  Rather than
separate AND/OR/XOR/NAND instructions, SM89 encodes an 8-bit truth table
in the instruction word itself.  Any Boolean function of three 32-bit inputs
can be computed in a single 2.51-cycle operation.  This is architecturally
elegant: the compiler emits \texttt{LOP3.LUT 0xc0} for AND, \texttt{0xfc}
for OR, \texttt{0x3c} for XOR, and arbitrary combinations for complex
predicates.  The uniform variant \texttt{ULOP3.LUT} operates on uniform
registers with the same truth-table encoding.

\paragraph{8.2}
\texttt{POPC} (population count / bit count) at 11.77 cycles per dependent
pair is notably slower than other integer operations.  This confirms it is
a multi-cycle integer operation, not an SFU operation as some documentation
suggests.  The per-pair measurement artifact (0.26cy/pair for zero-latency
operations like \texttt{IABS}) was carefully subtracted.

\textit{Source:} \texttt{probes/probe\_bitwise.cu},
\texttt{probes/probe\_bit\_manipulation.cu}.

%----------------------------------------------------------------------
\subsection{Type Conversion}
\label{par:cvt}

\paragraph{9.1}
Conversion instructions (\texttt{F2F}, \texttt{I2F}, \texttt{F2I},
\texttt{I2I}) uniformly measure at 4.54 cycles, identical to FP32 FMA.
This places them in the FMA pipeline, not a separate conversion unit.
A surprising finding: \texttt{FABS} and \texttt{FNEG} measure as
\emph{free} -- zero additional latency when chained with arithmetic.
These are pipeline modifiers (sign-bit manipulation in the register
read port), not independent operations.  Similarly, \texttt{.SAT}
(saturation clamp to [0,1]) is free when used as an FFMA modifier.

\textit{Source:} \texttt{microbench/microbench\_latency\_conversions.cu},
\texttt{probes/probe\_conversions.cu}.

%----------------------------------------------------------------------
\subsection{Global Memory Subsystem}
\label{par:memory}

\paragraph{10.1}
\texttt{LDG} (global load) at 92.29 cycles via pointer-chasing measures
the full L1 miss + L2 hit latency path.  This is the dominant cost in
bandwidth-bound kernels.  The vectorized forms (\texttt{LDG.E.64},
\texttt{LDG.E.128}) do not reduce latency but increase bandwidth per
instruction, which is why vectorized loads are critical for memory-bound
code.  The \texttt{.CONSTANT} modifier routes loads through the constant
cache (LDC path) which has different latency characteristics.

\paragraph{10.2}
The cache policy modifiers (\texttt{.L1}, \texttt{.L2},
\texttt{.STRONG.GPU}, \texttt{.STRONG.SYS}) create qualitatively different
execution regimes.  The controlled cache-policy experiments showed that
\texttt{-dlcm=cg} (forcing \texttt{STRONG.GPU}) changes emitted load
spellings but does not materially change cycle counts or stall
distributions once the SYS64 fused body exists.  Cache policy is a
secondary refinement, not a primary driver.

\textit{Source:} \texttt{probes/probe\_memory.cu},
\texttt{probes/data\_movement/}, \texttt{probes/cache\_control/}.

%----------------------------------------------------------------------
\subsection{Shared Memory Subsystem}
\label{par:smem}

\paragraph{11.1}
\texttt{LDS} (shared memory load) at 28.03 cycles via pointer-chasing
measures on-chip SRAM latency.  This is 3.3x faster than global memory,
confirming the critical importance of shared memory tiling for
bandwidth-bound kernels.  Bank conflicts add latency multiplicatively --
the \texttt{microbench\_smem\_bank\_conflicts} kernel quantifies this at
up to 32x for fully serialized 32-way conflicts.

\textit{Source:} \texttt{probes/probe\_memory.cu},
\texttt{microbench/microbench\_smem\_bank\_conflicts.cu}.

%----------------------------------------------------------------------
\subsection{Asynchronous Copy Pipeline}
\label{par:async}

\paragraph{12.1}
\texttt{LDGSTS} (load-global-store-shared) is the async copy instruction
that bypasses the register file.  It is the backbone of the modern CUDA
pipeline: data moves from global memory directly to shared memory while
the SM continues executing arithmetic.  The associated \texttt{LDGDEPBAR}
and \texttt{DEPBAR.LE} instructions form the dependency tracking system
that tells the warp scheduler when the async copy has completed.

This trio (\texttt{LDGSTS + LDGDEPBAR + DEPBAR.LE}) is the ``anchor''
that the combo frontier research identified as the central enabling
structure for modern GPU kernel design.

\textit{Source:} \texttt{probes/data\_movement/probe\_cp\_async\_zfill.cu},
\texttt{probes/probe\_async\_copy.cu}.

%----------------------------------------------------------------------
\subsection{Atomic Operations}
\label{par:atomic}

\paragraph{13.1}
The atomic instruction family spans three scopes: \texttt{ATOMS} (shared
memory), \texttt{ATOM}/\texttt{ATOMG} (global memory), and system-scope
variants (\texttt{ATOMG.E.*.STRONG.SYS}).  The scope transition from
\texttt{.GPU} to \texttt{.SYS} introduces \texttt{MEMBAR} as a major
stall source, pushing the kernel into a qualitatively different scheduling
regime.  This was the most important architectural finding from the combo
frontier experiments.

\textit{Source:} \texttt{probes/atomic\_sweep/},
\texttt{probes/shared\_atomics/}, \texttt{probes/edge\_atomics/}.

%----------------------------------------------------------------------
\subsection{Synchronization and Barriers}
\label{par:sync}

\paragraph{14.1}
\texttt{MEMBAR.GPU} at 205 cycles and \texttt{MEMBAR.SYS} at 2583 cycles
are the most expensive operations measured.  The 12.6x ratio between GPU
and SYS scope reveals the PCIe round-trip cost of system-scope memory
ordering.  For kernels that require cross-GPU coherence (multi-GPU
workloads), this cost dominates all other instruction latencies by an
order of magnitude.

\paragraph{14.2}
\texttt{BSSY} and \texttt{BSYNC} (structured synchronization) appear in
optimized code; \texttt{WARPSYNC} appears only in debug (\texttt{-G})
builds.  Both implement \texttt{\_\_syncwarp()} semantics, but the
compiler's optimization level determines which SASS form is selected.
This is a form-selection decision, not an ISA-level distinction.

\textit{Source:} \texttt{probes/barrier\_sync2/},
\texttt{probes/probe\_convergence\_barriers.cu}.

%----------------------------------------------------------------------
\subsection{Warp-Level Operations}
\label{par:warp}

\paragraph{15.1}
\texttt{SHFL.BFLY} (warp shuffle, butterfly pattern) at 24.96 cycles
enables register-to-register communication between threads in a warp
without shared memory.  The shuffle family (\texttt{.IDX}, \texttt{.UP},
\texttt{.DOWN}, \texttt{.BFLY}) provides the primitives for warp-level
reductions, scans, and neighbor communication.  The volume rendering
kernel in \texttt{instant\_ngp/} uses \texttt{SHFL} for warp-level
neighbor sharing of accumulated opacity, avoiding shared memory entirely.

\textit{Source:} \texttt{probes/probe\_memory.cu} (SHFL section),
\texttt{probes/simd\_video/}.

%----------------------------------------------------------------------
\subsection{Special Registers and Scheduling}
\label{par:special}

\paragraph{16.1}
\texttt{S2R} (special register read) provides access to hardware state:
thread indices (\texttt{SR\_TID.*}), clock counters (\texttt{SR\_CLOCKLO}),
warp identifiers, and SM topology.  \texttt{CS2R} reads 64-bit cycle
counters.  These are the foundation of all microbenchmark timing:
the latency measurements throughout this monograph are derived from
\texttt{CS2R SR\_CLOCKLO} bracketing of dependent instruction chains.

\texttt{NANOSLEEP} is a cooperative scheduling hint that yields the warp's
time slice.  It has no direct performance benefit but allows the warp
scheduler to issue other warps during spin-wait loops, improving overall
SM utilization.

\textit{Source:} \texttt{probes/probe\_special\_regs.cu},
\texttt{probes/probe\_nanosleep\_scheduling.cu}.

%----------------------------------------------------------------------
\subsection{Uniform Register Operations}
\label{par:uniform}

\paragraph{17.1}
The uniform instruction family (\texttt{ULDC}, \texttt{UIADD3},
\texttt{ULOP3}, \texttt{USHF}, \texttt{UISETP}) operates on a separate
register file of warp-uniform values.  Getting \texttt{UISETP} to appear
in SASS requires that compare operands remain in uniform registers through
scheduling, not merely that they are warp-uniform in source semantics.
The compiler aggressively rewrites the uniform compare chain when the
surrounding code becomes complex, falling back to per-thread \texttt{ISETP}.

\paragraph{17.2}
\texttt{UPLOP3.LUT} is the frontier instruction discovered through
cubin-side structural analysis.  It requires a specific structural path
(\texttt{PLOP3 -> UPLOP3}), not the intuitively expected
\texttt{ULOP3 -> UPLOP3}.  This structural law, proven by differential
fuzzing and compute-sanitizer validation, demonstrates that the uniform
and predicate subsystems have asymmetric coupling in the SM89 backend.

\textit{Source:} \texttt{probes/probe\_uniform\_datapath.cu},
\texttt{probes/probe\_uplop3\_uniform\_predicates.cu},
detailed analysis in \texttt{SM89\_UNIFORM\_LOP3\_ANALYSIS.md}.

%----------------------------------------------------------------------
\subsection{Predicate Register Operations}
\label{par:predicate}

\paragraph{18.1}
\texttt{P2R} (predicate-to-register) and \texttt{R2P}
(register-to-predicate) transfer values between the 7 predicate registers
(P0--P6) and the general-purpose register file.  The byte-qualified forms
\texttt{P2R.B1}, \texttt{P2R.B2}, \texttt{P2R.B3} pack predicate bits into
specific byte lanes of a GPR.  These forms are valid on SM89 (confirmed by
cubin substitution) but are never selected by any tested compiler frontend
(nvcc, clang, Triton) or ptxas version (CUDA 11.8, 12.6, 13.1).

This is the clearest example of the form-selection gap: the ISA supports
the instruction, the hardware executes it, but the compiler's cost model
consistently prefers alternative lowerings.  The P2R.B* frontier remains
the most significant open question in the SM89 characterization.

\textit{Source:} \texttt{probes/probe\_p2r\_*.cu},
\texttt{probes/probe\_predicate\_*.cu},
detailed analysis in \texttt{SM89\_PREDICATE\_REGISTER\_ANALYSIS.md}.
